% Source: https://artofproblemsolving.com/community/c3h2586463_number_theory
% Stable identity: aops-topic-2586463-q01
Write the equation as $\left(3^x-2^y\right)\left(3^x+2^y\right)=5^z$. The second factor is positive, so the first is too, and both are powers of $5$:
\[3^x-2^y=5^s,\qquad3^x+2^y=5^t,\qquad s+t=z,\quad s<t.\]
Subtracting, $2^{y+1}=5^s\left(5^{t-s}-1\right)$, so $5^s\mid2^{y+1}$, which forces $s=0$. Hence
\[3^x-2^y=1\qquad\text{and}\qquad3^x+2^y=5^z.\]
\textbf{Solve $3^x=2^y+1$.} If $y=0$, $3^x=2$ is impossible. If $y=1$, then $x=1$. If $y\ge2$, then $3^x\equiv1\pmod4$, so $x$ is even, $x=2m$, and $\left(3^m-1\right)\left(3^m+1\right)=2^y$. Two powers of $2$ differing by $2$ must be $2$ and $4$, so $3^m=3$, $x=2$ and $y=3$.

\textbf{Check the second equation.} For $(x,y)=(1,1)$: $3+2=5=5^1$, so $z=1$. For $(x,y)=(2,3)$: $9+8=17$ is not a power of $5$.

The only solution is $\boxed{(x,y,z)=(1,1,1)}$: $9-4=5$.
